\chapter{绪论}
\label{chap:introduction}

\section{研究背景}
\label{sec:introduction-research-background}

建筑幕墙是现代建筑的重要外围护结构，承担围护、采光、隔热、防水、防风压、节能和外观表达等多重功能。随着城市中大量公共建筑、商业建筑和高层建筑进入长期服役阶段，幕墙系统的安全性、耐久性和可维护性逐渐成为建筑运维中的重要问题。幕墙材料长期暴露在风、雨、温差、紫外线、污染物和偶发冲击环境中，可能出现玻璃爆裂、玻璃平整度异常、金属锈蚀、石材裂缝、石材污渍以及连接构件退化等问题。这些问题如果不能及时发现和评估，不仅会影响建筑外观，还可能对围护性能、使用安全和维护成本产生持续影响。

在建筑运维数字化和 BIM\footnote{BIM（Building Information Modeling），即建筑信息模型。它通过数字化方式组织建筑的几何形态、构件属性、空间关系和运维信息，使建筑全生命周期中的设计、施工、检测和维护数据能够被统一表达和管理\cite{gbt51212-2016}。} 技术发展的背景下，幕墙评估的目标正在从“发现某个缺陷”转向“理解建筑外围护系统的整体状态”。韧性概念强调系统在扰动或退化条件下维持关键功能、恢复功能和适应环境变化的能力\cite{shi2023resilientcity}。对于建筑幕墙而言，韧性评估不仅需要知道某张图像中是否存在缺陷，还需要回答缺陷分布在哪些区域、不同材料上的损伤类型是否具有聚集趋势、局部损伤对整体运维决策有何意义、后续维护应优先关注哪些问题等更高层次的问题。

现实工程中，幕墙巡检仍然大量依赖人工经验。检测人员需要面对高空作业、大面积立面、海量图像、复杂材料和不一致的现场环境。即便无人机和高清相机降低了图像采集难度，后续的图像整理、材质判断、缺陷识别、结果汇总和报告撰写仍然消耗大量人工。实际业务中经常出现“图像采集已经完成，但评估报告迟迟难以形成”的情况：现场已经获得大量巡检图像，但图像缺少面向建筑区域的分类整理，单张图像的检测结果彼此割裂，缺陷类型、分布区域、风险倾向和维护优先级需要工程人员反复比对、筛选和汇总。已有计算机视觉方法能够完成裂缝、锈蚀、破损等局部缺陷检测，但多数方案仍停留在单模型、单图像、单缺陷识别层面，只能检测某张图像是否存在某类缺陷，缺少多组巡检图像、多种幕墙材质、多类检测结果之间的汇总融合。如\cref{fig:introduction-traditional-inspection-and-user-demand} 所示，对于土木工程领域用户而言，实际需求并不是孤立的单图检测结论，而是面向整栋建筑形成的包括幕墙整体状态、主要缺陷类型、缺陷集中区域、风险倾向等多维度韧性评估结论。

\begin{figure}[!htb]
  \centering
  \resizebox{\linewidth}{!}{%
    \begin{tikzpicture}[
      font=\scriptsize,
      >=stealth,
      panel/.style={draw, rounded corners=2pt, line width=0.45pt, fill=gray!4},
      title/.style={font=\scriptsize\bfseries, align=center},
      flowstep/.style={draw, rounded corners=2pt, align=center, minimum width=3.65cm, minimum height=0.72cm, line width=0.45pt, inner sep=3pt},
      oldstep/.style={flowstep, fill=red!6},
      newstep/.style={flowstep, fill=blue!6},
      keyold/.style={flowstep, fill=red!12},
      keynew/.style={flowstep, fill=green!10},
      arrow/.style={->, line width=0.5pt},
      note/.style={align=center, font=\scriptsize, text=black}
    ]
      \draw[panel] (0, 0) rectangle (5.1, -6.9);
      \node[title] at (2.55, -0.45) {传统幕墙巡检流程};
      \node[oldstep] (old-collect) at (2.55, -1.35) {无人机或人工采集图像};
      \node[oldstep] (old-organize) at (2.55, -2.55) {人工整理图像文件};
      \node[oldstep] (old-detect) at (2.55, -3.75) {单图单缺陷检测};
      \node[oldstep] (old-merge) at (2.55, -4.95) {人工比对汇总结果};
      \node[keyold] (old-report) at (2.55, -6.15) {评估报告形成滞后};

      \draw[arrow] (old-collect) -- (old-organize);
      \draw[arrow] (old-organize) -- (old-detect);
      \draw[arrow] (old-detect) -- (old-merge);
      \draw[arrow] (old-merge) -- (old-report);

      \node[note] at (6.2, -3.45) {从单图割裂\\到建筑级融合};
      \draw[arrow] (5.35, -3.45) -- (7.05, -3.45);

      \draw[panel] (7.3, 0) rectangle (12.4, -6.9);
      \node[title] at (9.85, -0.45) {用户真实评估需求};
      \node[newstep] (new-assets) at (9.85, -1.35) {项目级图像统一管理};
      \node[newstep] (new-agent) at (9.85, -2.55) {按图像内容自动分流};
      \node[newstep] (new-detect) at (9.85, -3.75) {多类缺陷协同检测};
      \node[newstep] (new-summary) at (9.85, -4.95) {单图结果可读总结};
      \node[keynew] (new-report) at (9.85, -6.15) {建筑级韧性评估报告};

      \draw[arrow] (new-assets) -- (new-agent);
      \draw[arrow] (new-agent) -- (new-detect);
      \draw[arrow] (new-detect) -- (new-summary);
      \draw[arrow] (new-summary) -- (new-report);
    \end{tikzpicture}%
  }
  \caption{传统幕墙巡检流程与用户真实评估需求对比}
  \label{fig:introduction-traditional-inspection-and-user-demand}
\end{figure}

因此，本文关注的核心矛盾并不仅是单个检测模型能否识别某类缺陷，而是如何把图像采集、资产组织、任务调度、检测结果解释和建筑级报告生成连接为完整闭环。

\section{研究问题}
\label{sec:introduction-research-problem}

本文关注的研究问题可以概括为三个层次。

第一个层次是数据规模与组织问题。一个建筑项目可能包含大量幕墙图像，图像之间还需要按照立面、楼层、区域或采集批次进行组织。若缺少图像资产管理和分类整理能力，图像只能作为普通文件存在，难以支撑后续检索、复核、状态追踪和报告生成。更重要的是，建筑幕墙评估不能只停留在“图像数量很多”的层面，还需要从建筑整体维度理解不同区域的缺陷类型、空间分布、风险倾向和维护优先级。例如，同样是锈蚀缺陷，如果集中出现在某一立面或特定高度区间，其工程意义与零散出现完全不同；同样是污渍问题，如果与局部渗水或排水异常相关，也需要在报告中形成区域化判断。因此，系统需要同时支持项目级图像资产管理和建筑维度的区域化汇总能力。

第二个层次是任务选择与调度问题。幕墙图像中可能包含玻璃、金属、石材或混凝土等不同材料，不同材料对应的检测能力不同。传统固定流程会对每张图像执行全部检测模型，这会造成计算浪费和语义噪声。例如，对玻璃幕墙图像执行石材裂缝检测并无实际意义，对石材图像执行玻璃平整度检测也难以产生可信结论。另一种做法是由人工先判断图像材质，再手动把图像分配到不同检测模型中，但在批量巡检场景下，这种流程需要大量重复性人工分类，效率低、成本高，并且不同人员之间判断标准不一定一致。因此，系统需要一种能够根据图像内容自主选择检测能力的机制，将开放视觉语义判断转化为稳定的工程任务代码。

第三个层次是韧性评估与报告生成问题。检测模型通常输出区域坐标、像素占比、置信度、掩码图像和浮层图像等底层结果。这些结果对算法复核有价值，但并不能直接构成建筑运维人员可读的韧性评估报告。如果直接把所有检测结果拼接给大语言模型，海量数据内容会造成上下文过长、信息重复、重点不清和报告生成不稳定等问题。因此，系统需要一种从原子检测结果到图像级总结，再到项目级报告的中间组织机制，使单图单点检测结果能够逐步转化为面向建筑整体的缺陷概述、风险判断和维护建议。

\begin{figure}[!htb]
  \centering
  \resizebox{\linewidth}{!}{%
    \begin{tikzpicture}[
      font=\scriptsize,
      >=stealth,
      problem/.style={draw, rounded corners=2pt, align=center, minimum width=4.0cm, minimum height=0.82cm, line width=0.45pt, fill=red!6, inner sep=3pt},
      ability/.style={draw, rounded corners=2pt, align=center, minimum width=5.1cm, minimum height=0.82cm, line width=0.45pt, fill=blue!6, inner sep=3pt},
      group/.style={draw, rounded corners=2pt, line width=0.45pt, fill=gray!4},
      layer/.style={font=\scriptsize\bfseries, align=center},
      arrow/.style={->, line width=0.5pt}
    ]
      \path[use as bounding box] (-0.60, 0.85) rectangle (12.30, -3.62);
      \draw[group] (0.00, 0.75) rectangle (4.90, -3.52);
      \draw[group] (5.70, 0.75) rectangle (11.70, -3.52);

      \node[layer] at (2.45, 0.30) {研究问题层次};
      \node[layer] at (8.70, 0.30) {系统能力指向};

      \node[problem] (p-data) at (2.45, -0.56) {数据规模与组织问题};
      \node[ability] (a-data) at (8.70, -0.56) {项目 / 图像组 / 图像资产管理};

      \node[problem] (p-task) at (2.45, -1.61) {任务选择与调度问题};
      \node[ability] (a-task) at (8.70, -1.61) {材质识别、检测能力选择、任务代码生成};

      \node[problem] (p-report) at (2.45, -2.66) {韧性评估与报告生成问题};
      \node[ability] (a-report) at (8.70, -2.66) {风险判断、项目整体评估、维护建议};

      \draw[arrow] (p-data) -- (a-data);
      \draw[arrow] (p-task) -- (a-task);
      \draw[arrow] (p-report) -- (a-report);
    \end{tikzpicture}%
  }
  \caption{建筑幕墙韧性评估的三层研究问题模型}
  \label{fig:introduction-three-layer-research-problem-model}
\end{figure}

如\cref{fig:introduction-three-layer-research-problem-model} 所示，上述三个层次共同构成本文的研究问题：如何设计一套面向建筑幕墙韧性评估的软件系统，使其能够管理项目级图像资产，基于图像内容自主选择检测能力，调度原子检测工具，并将多源检测结果逐层整合为项目级评估报告。

\section{Agent 范式的场景适用性}
\label{sec:introduction-agent-paradigm-scenario-applicability}

传统软件系统擅长执行确定流程，但面对“先判断场景，再选择工具，再综合输出结论”的任务时，往往需要大量规则和人工配置。建筑幕墙检测正属于这种情况。图像是否属于幕墙、主要材料是什么、应该运行哪些检测能力、检测结果如何转化为工程语言，这些问题具有一定语义判断和上下文理解需求。单纯依靠硬编码规则难以覆盖复杂场景，单纯依靠检测模型又无法完成任务规划和报告表达。

AI Agent 范式为该问题提供了新的解决思路。近年来，学术界和工业界逐渐将大语言模型从“文本生成器”扩展为“任务规划和工具调用中心”。ReAct 方法强调推理和行动交替进行，使模型能够在推理过程中调用工具并利用工具反馈更新计划\cite{yao2023react}；Toolformer 讨论语言模型学习使用外部 API 的能力\cite{schick2023toolformer}；AutoGen 等框架进一步展示了多 Agent 协作和可编程智能工作流的工程价值\cite{wu2023autogen}。这些研究表明，Agent 的核心价值不只是输出自然语言，而是把理解、规划、工具调用和结果整合连接起来。如\cref{fig:introduction-ai-agent-paradigm-workflow} 所示，AI Agent 范式通常不是一次性生成答案，而是在任务输入、语义理解、任务规划、工具调用、观察反馈和结构化输出之间形成闭环。

\begin{figure}[!htb]
  \centering
  \resizebox{\linewidth}{!}{%
    \begin{tikzpicture}[
      font=\scriptsize,
      >=stealth,
      node/.style={draw, rounded corners=2pt, align=center, minimum width=2.18cm, minimum height=0.82cm, line width=0.45pt, inner sep=3pt},
      input/.style={node, fill=blue!6},
      agent/.style={node, fill=green!9},
      tool/.style={node, fill=orange!10},
      output/.style={node, fill=purple!7},
      panel/.style={draw, rounded corners=2pt, line width=0.45pt, fill=gray!4},
      arrow/.style={->, line width=0.5pt},
      loop/.style={->, dashed, line width=0.45pt}
    ]
      \draw[panel] (-0.2, 0.55) rectangle (12.5, -3.75);

      \node[input] (task) at (1.15, -0.73) {任务输入\\与上下文};
      \node[agent] (understand) at (3.65, -0.73) {语义理解\\与目标识别};
      \node[agent] (plan) at (6.15, -0.73) {任务规划\\与工具选择};
      \node[tool] (execute) at (8.65, -0.73) {工具调用\\与外部执行};
      \node[output] (result) at (11.15, -0.73) {结构化输出\\与状态反馈};
      \node[agent] (observe) at (8.65, -2.48) {观察反馈\\与结果校验};

      \draw[arrow] (task) -- (understand);
      \draw[arrow] (understand) -- (plan);
      \draw[arrow] (plan) -- (execute);
      \draw[arrow] (execute) -- (result);
      \draw[arrow] (execute) -- (observe);
      \draw[arrow] (observe) -- (result);
      \draw[loop] (observe.west) to[out=180, in=-90] (plan.south);
    \end{tikzpicture}%
  }
  \caption{AI Agent 范式的一般工作流}
  \label{fig:introduction-ai-agent-paradigm-workflow}
\end{figure}

本文将 Agent 引入建筑幕墙韧性评估，并不是为了给系统增加一个对话入口，而是为了让 Agent 承担三个具体职责：在检测前完成任务路由，在检测后完成图像级语义压缩，在项目结束时完成报告级信息融合。这样的设计使 Agent 真正嵌入业务链路，而不是作为孤立的问答功能存在。

如果采用传统规则系统，也可以通过人工配置材料和检测任务之间的映射关系。例如，当用户选择“玻璃幕墙”时，系统固定执行玻璃平整度和玻璃爆裂检测；当用户选择“石材幕墙”时，系统固定执行裂缝和污渍检测。但是现实图像往往并不如此简单。一张图像可能同时包含玻璃和金属框架，也可能包含非幕墙背景、反射、遮挡或局部构件。完全依赖用户手动选择材料，会降低批量处理效率；完全依赖硬编码规则，又难以处理复杂视觉语义。Agent 的价值正在于它可以在有限工具集合内完成语义判断，从而把复杂场景判断转化为工程系统可执行的任务代码。

如果直接让大语言模型生成最终报告，也存在明显问题。幕墙项目可能包含多张图像，每张图像又包含多个检测任务和多个区域结果。原始结果中大量坐标、像素、置信度和图像产物信息会干扰报告生成。本文没有让 Agent 一步完成所有任务，而是将其拆分为三个受控节点：分类路由、图像级总结和项目报告。这种设计既符合当前 Agent 工程化中的工具调用思想，也符合建筑韧性评估从局部观察到整体判断的认知过程。

\section{系统技术定位}
\label{sec:introduction-system-technical-positioning}

本项目是一个以工程系统为载体、以 Agent 能力为核心亮点的毕业设计。工程系统本身采用的是相对成熟的技术组合，例如前后端分离、gRPC、消息队列、对象存储、数据库事务和容器化部署等。这些工程技术是系统稳定性、可追踪性和可部署性的基础。没有工程链路，Agent 只能处理单次输入；有了工程链路，Agent 的判断才能驱动任务创建、状态推进、结果落库、异常恢复和报告生成。

因此，本系统的技术定位有别于“实现一个视觉检测模型”或“实现一个通用后台管理系统”，而是“构建一个 Agent 驱动的幕墙缺陷检测与韧性评估方法框架”。工程链路提供高可用、高性能和高并发的运行保障，视觉检测能力提供原子事实来源，Agent 自主任务调度方法和 LLM 分层信息整合策略则实现真正的 AI 赋能。换言之，工程链路实现保证系统可用性，Agent 方法决定系统能否从单点检测提升到建筑维度评估。

\section{研究目标}
\label{sec:introduction-research-goal}

本文的研究目标是设计并实现一套基于 AI Agent 的建筑幕墙韧性评估软件系统。系统需要完成从项目创建、图像资产管理、Agent 智能检测、人工复核确认到评估报告生成的完整流程。在智能检测阶段，本文重点研究面向建筑幕墙多材质图像的 Agent 自主任务调度方法，通过分类 Agent 判断每张图像需要执行哪些原子检测能力；在原子检测阶段，系统需要稳定调度金属锈蚀、石材裂缝、石材污渍、玻璃平整度和玻璃爆裂五类检测能力；在信息整合阶段，本文重点研究 LLM 分层信息整合策略，将单张图像的检测结果生成图像级总结，再将所有图像总结、项目背景、图像分组和复核信息汇总为项目级韧性评估报告。

从工程目标看，系统需要具备清晰模块边界和可部署性。前端需要提供完整项目流程和结果展示，后端需要支持长耗时任务、异步回调、实时状态推送、数据库事务、对象存储和容器化部署。工程实现不是本文的唯一亮点，但它是 Agent 能力落地的必要保障。没有稳定工程链路，Agent 自主任务调度方法和 LLM 分层信息整合策略只能停留在脚本演示阶段，无法成为实际可用的软件系统。

\section{研究内容与贡献}
\label{sec:introduction-research-content-and-contribution}

本文主要研究内容包括：

\begin{enumerate}
  \item 调研建筑幕墙检测、BIM 与 AI Agent 应用范式，分析现有单点检测方案在建筑维度评估中的不足；
  \item 建立幕墙韧性评估软件的问题模型，明确图像资产、检测任务、人工复核和报告生成之间的关系；
  \item 提出面向建筑幕墙多材质图像的 Agent 自主任务调度方法，使系统能够根据图像内容动态选择检测能力；
  \item 提出 LLM 分层信息整合策略，将原子检测结果先转化为图像级总结，再融合为项目级韧性评估报告；
  \item 交付可运行软件原型系统，实现前端交互、服务通信、异步任务、状态推送、对象存储、容器化部署和实际场景试用。
\end{enumerate}

本文贡献主要体现在以下三个方面：

\begin{enumerate}
  \item 将 Agent 工具调用思想引入建筑幕墙评估场景，提出 Agent 自主任务调度方法，使 Agent 的语义判断能够被工程状态机消费，降低批量场景下人工材质分类和任务选择成本；
  \item 提出面向建筑幕墙多材质图像的 Agent 自主任务调度方法，使系统能够根据图像内容动态选择检测能力；
  \item 提出 LLM 分层信息整合策略，用图像级总结作为项目级报告生成的中间语义层，解决海量检测结果直接汇总的上下文压力和信息噪声问题，使单图检测结论能够上升为建筑维度评估结论；
  \item 构建完整的工程闭环，将 Agent 自主任务调度方法、视觉检测模型、结构化落库、人工复核和报告生成连接为可部署系统，体现了工程实践与 Agent 方法结合的价值。
\end{enumerate}

\section{论文组织结构}
\label{sec:introduction-thesis-organization}

全文共七章。第一章介绍研究背景与问题、Agent 范式的场景适用性、系统技术定位、研究目标、内容与贡献。第二章梳理建筑幕墙缺陷检测、计算机视觉方法、AI Agent 范式、建筑幕墙韧性评估的相关研究基础。第三章对问题进行建模，分析系统需求、角色、数据对象和相关约束条件。第四章详细阐述本文提出的 Agent 自主任务调度方法和 LLM 分层信息整合策略等技术方案设计。第五章介绍系统架构与工程实现，重点说明如何构建高可用、高性能和高并发的工程链路，并支撑 Agent 能力稳定落地。第六章进行系统验证与应用分析，说明系统在功能、流程、部署和试用层面的完成情况。第七章总结全文工作内容并展望后续研究方向。
